\section{Exact radial coarea and Borel critical-origin bands}
\label{sec:rc-main}
The preceding critical-sector theorem takes a local limit at fixed
collision count and requires a tangent sector for each physical cut.
Here we replace the local Taylor calculation, for an explicitly defined
physical class, by an exact nonlinear calculation. This gives a uniform
source estimate at every band width, including arbitrary Borel cuts. It
neither identifies different physical roof levels nor bounds the paired
endpoint measure by the collision measure without a comparison proof.

Throughout this section the dynamics, the probability $\nu_R^*$, the
section, and the exact event $\Pi_{n,k,m,R}$ are those of the original
model. We write $\Phi_m(x)=(p_0(x),p_m(x))=\zeta$ and
$D=\operatorname{diag}(-1,1)$. The new polar radius $r=|\zeta|$ is
an auxiliary coordinate, not the physical clock and not one of the
microscopic widths $r_B$. Fix $0<\rho<1$ and put
\begin{equation}\label{eq:rc-constants}
 K_\rho=\frac{100}{47}\sqrt{1-\rho^2},\qquad
 U_m=\begin{cases}320/9,&m=1,\\170/9,&m\ge2.\end{cases}
\end{equation}
All estimates below are uniform in $R\in[9/20,47/100]$ and in $m$.
There is no cutoff on a positive interior incidence.

\begin{lemma}[An endpoint-momentum metric away from the critical point]
\label{lem:rc-metric}
On every regular physical $m$-flight branch with $|\Phi_m|<\rho$, the
reduced action and its Hessian satisfy
\begin{equation}\label{eq:rc-metric}
 \widehat G:=R^{-2}D_q^2F,\qquad
 K_\rho I\preceq\widehat G\preceq U_m I.
\end{equation}
The map from endpoint angles $q=(\alpha_0,\alpha_m)$ to $\zeta$ is a
local diffeomorphism. For every local inverse $q=q(\zeta)$,
\begin{equation}\label{eq:rc-gradient}
 D_\zeta q=R G^{-1}D,\qquad
 \nabla_\zeta(F\circ q)=D\widehat G^{-1}D\zeta.
\end{equation}
These are identities for the actual finite itinerary on its regular
physical side; they do not define a physical collision map at tangency.
\end{lemma}
\begin{proof}
The full second variation and its minimizing Schur complement were
proved in Lemma~\ref{lem:v109-hessian}. They give
$G\succeq R\operatorname{diag}(c_0,c_m)$. Since
$|p_0|,|p_m|<\rho$, division by $R^2$ gives the lower bound in
\eqref{eq:rc-metric}. For $m\ge2$ the unreduced endpoint block is
 diagonal. Its two entries, divided by $R^2$, are
$c_0/R+c_0^2/\tau_0$ and $c_m/R+c_m^2/\tau_{m-1}$.
Each is at most $20/9+50/3=170/9$. Minimizing over the interior
variations subtracts a positive semidefinite matrix. For $m=1$ there
is no interior block. The absolute mixed entry, divided by $R^2$,
is $c_0c_m/\tau_0$; the row-sum bound is at most
$20/9+100/3=320/9$. Thus the upper estimate, previously used only at
normal endpoints in Lemma~\ref{lem:mc-hessian}, holds everywhere on
the stated regular physical set.

The exact first variation is $\nabla_qF=RD\zeta$, and hence
$D_q\zeta=R^{-1}DG$. Its determinant is nonzero by positivity of
$G$. Inverting and applying the chain rule gives
\[
 (D_\zeta q)^{\mathsf T}\nabla_qF
 =R D G^{-1}(RD\zeta)=D\widehat G^{-1}D\zeta.
\]
In particular, no derivative of a Hessian-dependent angular change of
variables has been dropped. Positivity is used only at regular
points with the true specular stationarity equations satisfied.
\end{proof}

\begin{theorem}[Exact radial roof disintegration and synchronized bound]
\label{thm:rc-coarea}
Let $e_\theta=(\cos\theta,\sin\theta)$ and, on any regular inverse
branch $x=x(\zeta)$, set
\begin{equation}\label{eq:rc-h-and-u}
 h_x(\zeta)=e_\theta^{\mathsf T}D\widehat G_x^{-1}De_\theta,
 \quad
 u_x(\zeta)=|\partial_{\alpha_0}p_m(x)|^{-1},
 \quad \zeta=re_\theta\ne0.
\end{equation}
For fixed $\theta$, while the same actual branch is regular,
\begin{equation}\label{eq:rc-radial}
 \frac{dL(x(re_\theta))}{dr}=r h_x(re_\theta),\qquad
 U_m^{-1}\le h_x\le K_\rho^{-1}.
\end{equation}
The collision-source density in $(L,\theta)$ on this branch is exactly
\begin{equation}\label{eq:rc-coarea-density}
 b_x(re_\theta)=\frac{u_x(re_\theta)}{4\pi c_*h_x(re_\theta)}.
\end{equation}
A Borel selector $0\le w\le1$, including the original initial/terminal
section tests, exact labels, roof intervals, guards and pairing cuts,
multiplies this density without a smoothness assumption. For the
complete one-copy regular physical root family,
\begin{equation}\label{eq:rc-synchronized}
 \sum_{x:\,\Phi_m(x)=re_\theta} w(x)b_x(re_\theta)
       \le \frac{U_m C_{\rm mom}}{4\pi c_*}
\end{equation}
for $0<r<\rho$ at every regular root value (and for the corresponding
almost-everywhere versions of measurable disintegrations).
The summands in \eqref{eq:rc-synchronized} share the same endpoint
momenta; their physical roofs need not be equal.
\end{theorem}
\begin{proof}
Equation~\eqref{eq:rc-gradient} gives \eqref{eq:rc-radial}. The
eigenvalue bounds in \eqref{eq:rc-metric} give its two-sided constants.
The determinant of the map
$(\alpha_0,p_0)\mapsto(p_0,p_m)$ is
$-\partial_{\alpha_0}p_m$. This derivative is nonzero: equivalently,
if $G_{0m}=e\ne0$, the generating matrix has lower-left entry
$-\det G/(Re)$. The nonvanishing of $e$ is the exact positive optical
denominator identity, not an assumed genericity condition.
Consequently the original probability contributes
$u_x\,d\zeta/(4\pi c_*)$, with the initial-section indicator retained.
Polar coordinates give $d\zeta=r\,dr\,d\theta$. At fixed angle,
$dt=r h_x\,dr$, where $t=L(x)$. The factors $r$ cancel, proving
\eqref{eq:rc-coarea-density}, with no missing $R$ or incidence factor.

For precision with arbitrary cuts, cover the regular source by
countably many relatively compact inverse charts, fix one angular
seam, and disjointize their source overlaps. Apply change of variables
first to bounded nonnegative tests on finitely many charts and then
use monotone convergence. On each angular slice, subdivide the open
regular inverse domain into intervals; on every such interval
\eqref{eq:rc-radial} is strictly positive. Borel cuts stay inside the
integrands rather than being turned into new open branches. This
proves the exact formula also when a cut has no angular density at
its boundary. The angular seam is only a coordinate seam and is
counted once. This argument does not discard a singular endpoint trace.

Finally $1/h_x\le U_m$, and Theorem~\ref{thm:mc-kernel} bounds the
complete root sum $\sum_xu_x(\zeta)$ by $C_{\rm mom}$. Apply it to
finite subfamilies and exhaust the nonnegative terms. Multiplying by
selectors between zero and one cannot increase this sum. The root
bound at regular root values was proved in that theorem; no chosen
version of an almost-everywhere density is evaluated at a new point.
\end{proof}

\subsection{Critical-origin arms with their physical cuts still present}
The next class is defined by an actual continuation property, not by
deleting an inconvenient source. A \emph{critical-origin radial arm}
is a regular physical inverse along a fixed momentum ray, defined for
$0<r<r_+(\theta)\le\rho$, with no missing collision on that interval.
The arm stays on the same strict visited word. Its endpoint at $r=0$
may be a regular critical orbit, an unvisited/section boundary, or a
limit with a visited tangency; no physical derivative is asserted at
that endpoint. All physical points with $r>0$ remain genuine regular
orbits. Only one copy of each physical initial point may be used.

We work with measurable one-copy families of these arms and arbitrary
Borel selectors on them. In particular this includes every original
pairing sheet whose incoming limit has $|\nabla_qF|=0$ and whose
regular interval reaches that limit. The original sheet construction
and its measurable, one-copy parametrization remain the explicitly
retained input of Section~\ref{sec:v109-inputs}; the present theorem
does not purport to reconstruct that historical construction. Its
geometric estimate applies to every family satisfying the definition,
not only to finitely many coordinate balls or to smooth selectors.
Other components of the original event, including arms which terminate
at positive momentum radius before reaching zero, are not assigned
zero source and remain outside this particular estimate.

\begin{lemma}[Exact band geometry, without a Taylor radius]
\label{lem:rc-band-geometry}
On a critical-origin arm the finite limit
$a(\theta)=\lim_{r\downarrow0}L(x(re_\theta))$ exists and
\begin{equation}\label{eq:rc-band-geometry}
 \frac{r^2}{2U_m}\le L(x(re_\theta))-a(\theta)
                      \le\frac{r^2}{2K_\rho}.
\end{equation}
No local asymptotic error occurs. The limiting endpoint angles also
exist in local angular lifts. The assertions remain valid when an
interior incidence tends to zero as $r\downarrow0$.
\end{lemma}
\begin{proof}
For $0<\eta<r$, integrate \eqref{eq:rc-radial} to obtain
\[
 \frac{r^2-\eta^2}{2U_m}\le L(r)-L(\eta)
                      \le\frac{r^2-\eta^2}{2K_\rho}.
\]
Thus $L(\eta)$ is Cauchy as $\eta\downarrow0$ and has a finite limit.
Taking this limit proves \eqref{eq:rc-band-geometry}. Also
$|dq/dr|\le1/(R K_\rho)$ by \eqref{eq:rc-gradient}; hence both
endpoint angles are Cauchy in consistent lifts on a sufficiently short
terminal interval. These estimates use no interior-incidence lower
bound. All calculations were made at $\eta>0$ on actual regular
orbits. Taking a limit of their physical lengths and endpoint angles
does not continue a collision map through a grazing collision.
\end{proof}

\begin{proposition}[No drift of the limiting roof on a physical fan]
\label{prop:rc-critical-roof-constancy}
Suppose a continued physical inverse has domain
$\{re_\theta:0<r<r_+(\theta),\ \theta\in I\}$, where $I$ is a
connected angular interval and every compact $J\subset I$ admits a
positive radius $r_J$ with $(0,r_J)\times J$ in that domain. Then
$a(\theta)$ in Lemma~\ref{lem:rc-band-geometry} is constant on $I$.
In particular the bands on such a fan are measured from one actual
limiting physical roof, not a freely chosen angular clock correction.
Arbitrary Borel cuts may still be imposed on the actual source after
this identity has been obtained.
\end{proposition}
\begin{proof}
Put $e_\theta^\perp=(-\sin\theta,\cos\theta)$. The exact angular
chain rule, using \eqref{eq:rc-gradient}, is
\[
 \partial_\theta L(re_\theta)
   =r^2(e_\theta^\perp)^{\mathsf T}
                         D\widehat G^{-1}D e_\theta .
\]
Its absolute value is at most $r^2/K_\rho$. Given two angles in $I$,
choose a compact connecting interval $J$ and use the common positive
radius in the hypothesis. Integration at any $0<r<r_J$ gives
$|L(re_\theta)-L(re_\psi)|\le r^2|\theta-\psi|/K_\rho$.
Letting $r\downarrow0$ proves equality of the two limiting roofs.
Only a compact subinterval is used; no collision-uniform continuation
radius is assumed. No angular boundary limit of a Borel selector is
needed because the identity is for the underlying physical action.
\end{proof}

\begin{theorem}[Complete Borel critical-band and singularity budgets]
\label{thm:rc-band}
For any measurable one-copy family of critical-origin arms let
$E_s$ be the set of its actual points satisfying
$0<L-a(\theta)\le s$, where $s>0$, and put
\[
 A_s=\min\{\rho,\sqrt{2U_m s}\}.
\]
For any Borel $0\le w\le1$, retaining the original event and original
probability, one has the simultaneous bounds
\begin{align}
 \int_{E_s}\one_{\Pi_{n,k,m,R}}w\,d\nu_R^*
   &\le\frac{C_{\rm mom}}{4c_*}A_s^2,
                                      \label{eq:rc-band-mass}\\
 \int_{E_s}\one_{\Pi_{n,k,m,R}}w\,|\Phi_m|^{-\gamma}\,d\nu_R^*
   &\le\frac{C_{\rm mom}}{2c_*(2-\gamma)}A_s^{2-\gamma},
     \qquad 0\le\gamma<2,             \label{eq:rc-band-weight}\\
 \int_{E_s}\one_{\Pi_{n,k,m,R}}w\,|\nabla_qF|^{-1}\,d\nu_R^*
   &\le\frac{C_{\rm mom}}{2c_*R}A_s.
                                      \label{eq:rc-inverse-gradient}
\end{align}
These estimates hold for the complete family before any supremum over
$R,m,n,k$, and for every $s$, not merely in a local-radius limit at a
fixed collision count. They also hold for a disjoint union of exact
records whose total selector is at most one. No tangent-sector limit,
finite $C^1$ description of a pairing cut, common chart radius, or
uniform bound on third derivatives of the roof is required.
\end{theorem}
\begin{proof}
By \eqref{eq:rc-band-geometry}, every selected point has
$|\Phi_m|\le A_s$. The positive joint-kernel domination from
Theorem~\ref{thm:mc-kernel}, on the \emph{actual} physical source, gives
\[
 \int_{E_s}\one_\Pi w\,|\Phi_m|^{-\gamma}\,d\nu_R^*
 \le\frac{C_{\rm mom}}{4\pi c_*}
       \int_{|\zeta|\le A_s}|\zeta|^{-\gamma}\,d\zeta.
\]
Apply this first with bounded truncations on finitely many disjoint
charts. Monotone convergence supplies the entire physical source and
the unbounded nonnegative weight. The polar integral is
$2\pi A_s^{2-\gamma}/(2-\gamma)$, proving
\eqref{eq:rc-band-weight}. Its case $\gamma=0$ is
\eqref{eq:rc-band-mass}. Finally $|\nabla_qF|=R|\Phi_m|$ gives
\eqref{eq:rc-inverse-gradient} at $\gamma=1$.

The constants never depended on how many arms were selected or on
their individual continuation lengths. A shorter actual arm contributes
only its actual points; it was not extended through an unvisited disk.
The same argument covers limiting visited tangencies because all
integrated points have positive visited incidences. It controls their
interior source approach, not an arbitrarily prescribed trace measure
supported on the singular boundary. General Borel cuts remain measurable
selectors throughout; their pointwise or angular boundary limits are
not assumed to exist. The assertion for unions follows from one-copy
positive domination, not from adding separately normalized probabilities.
\end{proof}

\subsection{What has and has not been transferred to the original roof}
Theorems~\ref{thm:rc-coarea} and \ref{thm:rc-band} remove the need
for Taylor remainders and tangent sectors for the stated synchronized
coarea and physical-band budgets. In particular
\eqref{eq:rc-band-mass} is a collision-uniform $O(s)$ statement and
\eqref{eq:rc-inverse-gradient} is a collision-uniform $O(\sqrt s)$
statement on the full defined arm family. These are continuum estimates,
not conclusions of a finite search or a fitted exponent.

They do not close \eqref{eq:uc-open-limit}. In the actual roof density
at a common $t$, formula \eqref{eq:rc-coarea-density} must be evaluated
at the branch-dependent radius $r_x(t,\theta)$. The synchronized sum
\eqref{eq:rc-synchronized} evaluates all branches at the \emph{same}
$r$ and is not a bound for that different sum. Likewise $L-a(\theta)$
is used only to specify a physical subset; it is not substituted for
$L$ in the exact record, Fourier reference, conditioning or LLT.

The band budget bounds the original collision-source measure. The
paired measure $c^\varepsilon\,du\,d\theta$ uses the minimum of its
two endpoint traces and is not automatically a submeasure of the
interior density $f^\varepsilon\,du\,d\theta$. Its transfer still
requires the retained positive-variation/comparison argument. The
inverse-gradient estimate does not include the unbounded insertion
$I_m$, and no such insertion is silently absorbed into its constant.
Critical-origin arms do not exhaust the positive-radius singular cuts,
noncritical folds or low-load source. These sources, the first-order
clearance contribution, nongrazing capacity, and the original spectral,
arithmetic and conditional-denominator interfaces retain their exact
proof obligations. No factor $m^2$ is attached to the new unscaled
bounds at a fixed outer width and declared harmless. The original
ordered limit and four-coordinate raw local inversion target remain
unchanged and are not asserted to be proved here.
